\subsection{Task 6: Giá và điểm đánh giá --- Phát hiện lựa chọn tốt}

\subsubsection{Thiết kế Idiom}

\begin{table}[H]
\centering
\renewcommand{\arraystretch}{1.1}
\begin{tabularx}{\textwidth}{|l|l|X|}
\hline
\textbf{Idiom} & \multicolumn{2}{l|}{\textbf{Scatter Plot}} \\ \hline
\textbf{What}
& \multicolumn{2}{X|}{Price: Quantitative (Key trục X)} \\ \cline{2-3}
& \multicolumn{2}{X|}{review\_scores\_rating: Quantitative (Key trục Y)} \\ \cline{2-3}
& \multicolumn{2}{X|}{good\_deal\_flag (derived): Categorical (Good Deal / Normal)} \\ \cline{2-3}
& \multicolumn{2}{X|}{number\_of\_reviews: Quantitative (Size -- proxy độ tin cậy)} \\ \hline
\textbf{How}
& \textbf{Encode} &
\textbf{Mark:} Point (Circle) \newline
\textbf{Channel:} \newline
\hspace*{1em} Pos X: Price (Quantitative) \newline
\hspace*{1em} Pos Y: review\_scores\_rating (Quantitative) \newline
\hspace*{1em} Hue Color: good\_deal\_flag (Categorical -- xanh=Good Deal, xám=Normal) \newline
\hspace*{1em} Size: number\_of\_reviews (Quantitative) \newline
Reference Lines: Dọc: MEDIAN(price); Ngang: Constant = 4.8 (ngưỡng Good Deal) \\ \cline{2-3}
& \textbf{Manipulate} &
Hover để xem giá, rating, room\_type, neighbourhood, number\_of\_reviews. \\ \cline{2-3}
& \textbf{Facet} &
Superimpose: Reference Lines đặt lên scatter plot. \\ \cline{2-3}
& \textbf{Reduce} &
Filter: number\_of\_reviews $\geq$ 5 (loại listing ít đánh giá, không đủ tin cậy). \\ \hline
\textbf{Why} & \multicolumn{2}{X|}{
\textbf{produce $\rightarrow$ discover $\rightarrow$ compare} \newline
Scatter plot tối ưu cho phân tích 2Q correlation. 4 góc phần tư từ reference lines giúp nhận diện Good Deal (góc trên trái: price $<$ median VÀ rating $\geq$ 4.8).
} \\ \hline
\textbf{Scale} & \multicolumn{2}{X|}{
Color: 2 (Good Deal / Normal). \newline
Items: listing với $\geq$ 5 reviews, price\_is\_outlier = False.
} \\ \hline
\end{tabularx}
\end{table}

\begin{figure}[H]
    \centering
    \safeincludegraphics[width=0.9\linewidth]{charts/domain_task6_chart2.png}
    \caption{Hình 6.2. Giá niêm yết và điểm đánh giá --- phát hiện lựa chọn tốt}
    \label{fig:domain-task6-chart2}
\end{figure}

\subsubsection{Đánh giá biểu đồ}

\textbf{1. Tính biểu đạt (Expressiveness):}

\begin{itemize}
    \item Thuộc tính: Giá niêm yết / Price (Q), Điểm đánh giá / Review Scores Rating (Q), Cờ ``Món hời'' / Good Deal Flag (C), Số lượng đánh giá / Number of Reviews (Q -- Size), Median Price (Q -- Reference Line dọc), Ngưỡng Rating 4.8 (Q -- Reference Line ngang), price\_is\_outlier (C -- filter), Number of Reviews $\geq$ 5 (Q -- filter).
    \item Channel:
    \begin{itemize}
        \item PosX (Price -- Q): vị trí ngang thể hiện giá $\rightarrow$ đúng.
        \item PosY (Rating -- Q): vị trí dọc thể hiện chất lượng $\rightarrow$ đúng. Rating cao hơn = vị trí cao hơn nhất quán với convention `up = good'.
        \item Hue Color: phân biệt Good Deal / Normal (C) $\rightarrow$ đúng. Chỉ 2 giá trị, discriminability tuyệt đối.
        \item Size: thể hiện Number of Reviews (Q) $\rightarrow$ đúng. Size phù hợp cho Q thứ cấp (proxy trust level).
    \end{itemize}
    \item Đánh giá mức độ quan trọng:
    \begin{itemize}
        \item Ưu tiên 1 -- Vị trí hai trục (PosX và PosY): PosX/PosY là cặp kênh quan trọng nhất trong scatter plot. PosX mã hóa Price (Quantitative) và PosY mã hóa Review Scores Rating (Quantitative) --- cả hai đều là dữ liệu định lượng liên tục. Theo Mackinlay, Position on a common scale là kênh hiệu quả nhất cho dữ liệu định lượng, cho phép người xem nhận biết chính xác vị trí từng listing theo hai chiều, đồng thời suy luận về mối quan hệ (hoặc sự vắng mặt của tương quan) giữa giá và rating. Reference lines (median giá dọc, average rating ngang) được superimpose lên hai trục này để tạo framework bốn góc phần tư --- phục vụ trực tiếp task discover và compare của biểu đồ.
        \item Ưu tiên 2 -- Màu sắc (Color Hue): Color Hue mã hóa Good Deal Flag (Categorical --- Good Deal / Normal) bằng hai màu tương phản (xanh lá / xám). Theo Mackinlay, Color Hue là kênh phù hợp nhất cho dữ liệu danh mục nhị phân. Màu xanh lá nổi bật trên nền xám giúp người xem lập tức xác định vùng ``Good Deal'' trong không gian scatter plot --- phục vụ trực tiếp task discover (phát hiện cơ hội). Kênh này hoạt động song song với reference lines: PosX/PosY xác định vị trí, Color Hue xác nhận nhãn --- cả hai cùng nhau tạo ra thông điệp rõ ràng mà không cần người xem tự tính toán.
        \item Ưu tiên 3 -- Kích thước điểm (Size / Area): Size mã hóa Number of Reviews (Quantitative) --- listing có nhiều review tương ứng với điểm lớn hơn. Theo Mackinlay, Area không phải kênh mạnh nhất cho dữ liệu định lượng về độ chính xác, nhưng đây là kênh phụ với mục tiêu encode độ tin cậy (trust level) thay vì giá trị chính xác: người xem không cần biết số review cụ thể, chỉ cần nhận biết listing nào ``đã được nhiều người kiểm chứng'' để củng cố độ tin tưởng vào Good Deal được phát hiện.
    \end{itemize}
    \item Kết luận: Scatter plot là idiom tối ưu để phân tích mối quan hệ giữa 2 biến định lượng với phân nhóm categorical.
\end{itemize}

\textbf{2. Tính hiệu quả (Effectiveness):}

\begin{itemize}
    \item Độ chính xác (Accuracy): PosX và PosY (position on common scale) --- channel chính xác nhất. Kênh Size (Area) --- accuracy thấp hơn, khó so sánh chính xác số reviews; đây là channel phụ nên chấp nhận được.
    \item Khả năng phân biệt (Discriminability): 2 màu (xanh/xám) rất dễ phân biệt. Reference lines tạo 4 góc phần tư rõ ràng, giúp người xem định vị ngay vùng `good deal' (góc trên-trái). Overplotting tại rating 4.8--5.0 là hạn chế --- có thể cải thiện bằng jitter hoặc transparency.
    \item Khả năng tách biệt (Separability): PosX, PosY, Color và Size tách biệt tốt.
\end{itemize}

\subsubsection{Phân tích biểu đồ}

\begin{itemize}
    \item Vùng góc trên-trái (price $<$ median, rating $\geq$ 4.8) là vùng `Good Deal' --- tập trung nhiều điểm xanh, chủ yếu thuộc Brooklyn và Queens.
    \item Manhattan có nhiều listing ở phần bên phải chart (giá cao) với rating không tương xứng --- không phải lựa chọn tối ưu về cost-efficiency.
    \item Listing `Good Deal' có Size lớn (nhiều reviews) là những nơi đã được kiểm chứng bởi nhiều khách --- độ tin cậy cao.
    \item Quan trọng: Không có tương quan dương rõ ràng giữa giá cao và rating cao --- nhiều listing Brooklyn/Queens giá thấp vẫn đạt rating 4.7--5.0, chứng minh giá không phải là proxy của chất lượng trên Airbnb.
    \item Khuyến nghị: Ưu tiên listing trong vùng `Good Deal' tại Brooklyn/Queens với \seqsplit{number\_of\_reviews} $\geq$ 20--30.
\end{itemize}
